\documentclass[11pt,a4paper]{article}
\usepackage[margin=22mm,headheight=15pt]{geometry}
\usepackage{fontspec}
\usepackage{polyglossia}
\setmainlanguage{russian}
\setotherlanguage{english}
\setmainfont{DejaVu Serif}[ItalicFont=DejaVu Serif,BoldItalicFont=DejaVu Serif Bold,ItalicFeatures={FakeSlant=0.18},BoldItalicFeatures={FakeSlant=0.18}]
\setsansfont{DejaVu Sans}
\setmonofont{DejaVu Sans Mono}[Scale=0.86]
\usepackage{amsmath,amsthm}
\usepackage{unicode-math}
\setmathfont{Latin Modern Math}
\usepackage{booktabs,longtable,array,enumitem,xcolor,fancyhdr,ragged2e,needspace}
\usepackage{xurl}
\usepackage[unicode,colorlinks=true,linkcolor=blue!45!black,urlcolor=blue!55!black]{hyperref}
\setlength{\emergencystretch}{3em}
\setlength{\parskip}{5pt}
\setlength{\parindent}{0pt}
\setlist{itemsep=3pt,topsep=5pt}
\pagestyle{fancy}
\fancyhf{}
\fancyhead[L]{\small Приложение: на стадии проверки}
\fancyhead[R]{\small 26 октября 2026}
\fancyfoot[C]{\thepage}
\newcommand{\Aut}{\operatorname{Aut}}
\providecommand{\square}{\mdlgwhtsquare}
\title{\fontsize{18}{23}\selectfont Гармонические узлы и группы Кремоны\protect\\Локально нильпотентные дифференцирования}
\author{Дополнительное исследовательское приложение\\На стадии проверки}
\date{26 октября 2026 года}
\begin{document}
\raggedright
\maketitle
\begin{abstract}
\textbf{Статус на 26 октября 2026 года: предварительные результаты,
на стадии проверки.} Проверки выполнены в рамках исследовательской работы
с ИИ; внешнее независимое рецензирование и литературный приоритет не заявляются.
Приложение не заменяет исходную статью и её сертификаты.


Исследованы следствия явных выпрямлений трилистника Шастри и
$H(3,5,7)$ из указанного пользователем репозитория. Получены полные
бесконечные семейства выпрямляемых параметризаций $H(3,4,c)$ и
$H(3,5,c)$, точный относительный критерий переноса через LND,
явные координаты ядра дифференцирования из гипотезы E и
самостоятельные доказательства минимальных степеней слайсов.
Общая выпрямляемость всех $H(a,b,c)$ не установлена.
Мотивы, K-теория и гомологическая алгебра не используются.
\end{abstract}

\section{Основание, статус и воспроизводимость}
Проверена текущая версия репозитория
\href{https://github.com/qft98romanov-hash/abhyankar-moh-trefoil/tree/3cf41be3fb3b1108e1ec91dd23dff7bfcf59569e}{\texttt{qft98romanov-hash/abhyankar-moh-trefoil}},
commit \texttt{3cf41be3fb3b1108e1ec91dd23dff7bfcf59569e}.
Это рабочее математическое досье; заявления о первенстве результатов
требуют отдельной литературной проверки. Общие классические механизмы
отделены от их применения к проверенным формулам.

Точные запуски \texttt{verify.py} и
\texttt{verify\_printed\_formulas.py} успешно подтвердили оба выпрямления
и все двенадцать напечатанных прямых и обратных компонент.
Использована SymPy 1.14.0 с точной рациональной арифметикой и соотношением
$a^2-a+1=0$; числа с плавающей точкой для сертификатов не применялись.
GAP и Macaulay2 в этом запуске не выполнялись. Проверка обратимости
использует факторизации, точную делимость и инъективность локализации;
она не подменяется одной проверкой якобиана.

Поле коэффициентов базовых выпрямлений равно
\[
K=\mathbb Q[a]/(a^2-a+1)\simeq\mathbb Q(\sqrt{-3}).
\]
Проверенные модели:
\[
\gamma_T=(t^3-3t,t^4-4t^2,t^5-10t),\qquad
\gamma_H=(C_3,C_5,C_7).
\]
Аффинная замена $M(x,y,z)=(x,y+2,z-5x)$ переводит $\gamma_T$
в $(C_3,C_4,C_5)$, где $C_n(t)=2T_n(t/2)$.
Таким образом, именно три исходные параметризации снабжены
удостоверенными выпрямителями. Средняя координата $\gamma_H$ имеет
линейный член $+5t$, а не $+4t$.

\Needspace{5\baselineskip}
\paragraph{Главные результаты продолжения.}
\begin{enumerate}
\item При фиксированных взаимно простых $a,b\geq2$ построены явные
полиномиальные сдвиги третьей координаты; достаточно не более
$\varphi(a)\varphi(b)/2$ представителей индекса $c$.
\item Все $H(3,4,c)$ при $\gcd(c,12)=1$ и все $H(3,5,c)$ при
$\gcd(c,15)=1$ выпрямляются над $K$. Это бесконечно много
параметризаций, но не бесконечно много различных вещественных типов узлов.
\item Для конкретного дифференцирования гипотезы E получены явные
полиномиальные координаты ядра: оно изоморфно $\mathbb C^{[4]}$.
Это следствие проверенного выпрямителя, не решение общей проблемы сокращения.
\item Минимальные полные степени слайсов связанных LND равны $7$ для
трилистника Шастри и $11$ для $H(3,5,7)$. Доказательства и отдельные
рациональные сертификаты не зависят от выпрямлений репозитория.
\item Критерий центра переноса усилен до необходимого и достаточного
внутри указанной относительной конструкции. Он равносилен наличию
слайса у соответствующего LND.
\item Для естественной модификации $H(3,7,8)$ доказана невозможность
выполнить критерий центра. Это ограничение метода, не доказательство
невыпрямляемости самой кривой.
\end{enumerate}

\section{Гармонические вложения: точные редукции и конечные семейства}

Во всём этом разделе $k$ --- поле характеристики нуль. Для согласования
с формулами репозитория используем моническую нормировку
\[
 C_n(t)=2T_n(t/2),\qquad C_0=2,\quad C_1=t,\quad
 C_{n+1}=tC_n-C_{n-1}.
\]
Она связана со стандартной нормировкой $T_n(\cos\theta)=\cos(n\theta)$
линейной заменой параметра и масштабированием координат. Поэтому вопрос
выпрямляемости одинаков для обеих нормировок. Обозначим
$\gamma_{a,b,c}=(C_a,C_b,C_c)$.

\paragraph{Что уже известно в литературе.}
В статье G.~Freudenburg и J.~Freudenburg,
\emph{Curves Defined by Chebyshev Polynomials},
\href{https://arxiv.org/abs/0902.3440}{arXiv:0902.3440},
предложение 4.1, установлено: если $a,b,c\geq2$, то $\gamma_{a,b,c}$
является замкнутым вложением тогда и только тогда, когда индексы попарно
взаимно просты. При наличии индекса $1$ вложение и его выпрямление
очевидны, даже без взаимной простоты остальных индексов.
Предложение 4.2 той же статьи даёт алгебраическую редукцию индексов;
следствия 4.1 и 4.2 устанавливают выпрямляемость всех вложений
$H(2,b,c)$ и $H(a,b,a+b)$ над $k$.

Работы P.-V.~Koseleff и D.~Pecker
\href{https://arxiv.org/abs/0909.3281}{\emph{Chebyshev diagrams for two-bridge knots}}
и \href{https://arxiv.org/abs/1203.4376}{\emph{Harmonic knots}}
дают топологическую классификацию семейств с первым индексом не больше
$4$. Эта классификация сама по себе не является классификацией
вложений по полиномиальным автоморфизмам над $\mathbb C$.

\paragraph{Явный полином, восстанавливающий параметр.}
Основные тождества имеют вид
\[
 C_mC_n=C_{m+n}+C_{|m-n|},\qquad C_m\circ C_n=C_{mn}
 \quad(m,n\geq1).
\]
Пусть $a$ нечётно и $a,b,c$ попарно взаимно просты. Существуют
положительные целые $r,s,q$ такие, что
\[
 |rb-sc|=1,\qquad rb+sc=qa.
\]
Действительно, после выбора положительного решения $rb-sc=1$ можно
заменить $(r,s)$ на $(r+cj,s+bj)$; второе условие сводится к разрешимой
конгруэнции с коэффициентом $2bc$, обратимым по модулю $a$.
Тогда
\[
 \boxed{\;C_r(C_b(t))C_s(C_c(t))-C_q(C_a(t))=t.\;}
\]
Это сертификат равенства $k[C_a,C_b,C_c]=k[t]$. Он удостоверяет
\emph{вложение}, но не утверждает, что восстанавливающий полином
$C_r(y)C_s(z)-C_q(x)$ является координатой кольца $k[x,y,z]$.

\paragraph{Элементарная редукция индексов.}
Если $c=ra+sb$, $r,s>0$, то
\[
 (x,y,z)\longmapsto(x,y,C_r(x)C_s(y)-z)
\]
--- полиномиальная инволюция, переводящая
$\gamma_{a,b,c}$ в $\gamma_{a,b,|ra-sb|}$.
В стандартной нормировке $T_n$ перед произведением нужен множитель $2$.
В HTML-тексте доказательства предложения 4.2 статьи Freudenburg--Freudenburg
этот множитель в записи инволюции пропущен; предшествующее тождество
напечатано с правильным множителем.

Сумма индексов при редукции строго убывает. После перестановок координат
процесс заканчивается либо индексом $1$, либо тройкой
\[
 2\leq A<B<C,\qquad C\notin\langle A,B\rangle,
 \qquad C\leq AB-A-B.
\]
Редукция обеспечивает конечный список кандидатов при фиксированной
исходной паре индексов. Она не утверждает, что различные конечные
тройки алгебраически неэквивалентны, и не даёт выпрямления всех таких троек.

\subsection{Алгебраические сдвиги третьего индекса}

Ниже приведён самостоятельный элементарный вывод из описания узлов
плоской чебышёвской кривой. Заявления о новизне этого следствия не делается.
Он полезен тем, что переводит определённые сравнения индексов в
\emph{конкретные полиномиальные сдвиги} третьей координаты.

\paragraph{Лемма о кольце узловой проекции.}
Пусть $a,b\geq2$, $\gcd(a,b)=1$ и
\[
 A=k[C_a(t),C_b(t)]\subset k[t],\qquad
 \delta=\frac{(a-1)(b-1)}2.
\]
Над алгебраическим замыканием плоская кривая
$(C_a(t),C_b(t))$ имеет ровно $\delta$ обыкновенных двойных точек
и не имеет других особенностей (предложение 3.1 статьи
Freudenburg--Freudenburg). Если $(s_\nu,t_\nu)$ --- две прообразные
точки каждого узла, то
\[
 \boxed{\quad A=\{f\in k[t]:f(s_\nu)=f(t_\nu)\text{ для всех }\nu\}.\quad}
\]
Равенства понимаются после расширения поля, если узлы не расщепляются над $k$.

\emph{Доказательство.}
Включение слева направо очевидно. Полиномы
$C_a^rC_b^s$, $r\geq0$, $0\leq s<a$, образуют базис $A$ над $k$:
порождаемость следует из соотношения $C_a(C_b)=C_b(C_a)$, монического
по $C_b$ степени $a$. Их ведущие степени $ra+sb$ попарно различны.
Удаление ведущих членов показывает, что коразмерность $A$ в $k[t]$
равна числу пробелов полугруппы $\langle a,b\rangle$, то есть $\delta$.
С другой стороны, $\delta$ равенств значений независимы: все $2\delta$
прообразных точек узлов различны, и на них применима обычная полиномиальная
интерполяция. Поэтому правая часть тоже имеет коразмерность $\delta$.
Это доказывает равенство над алгебраическим замыканием, а алгоритм
удаления ведущих членов даёт спуск к $k$.\hfill$\square$

\paragraph{Теорема о сдвиге.}
Пусть целые $c,c'\geq1$ удовлетворяют
\[
 c'\equiv\varepsilon_a c\pmod {2a},\qquad
 c'\equiv\varepsilon_b c\pmod {2b},\qquad
 \varepsilon_a,\varepsilon_b\in\{1,-1\}.
\]
Положим $\sigma=\varepsilon_a\varepsilon_b$. Тогда существует
$P\in\mathbb Q[x,y]$ такое, что
\[
 \boxed{\quad C_{c'}(t)-\sigma C_c(t)=P(C_a(t),C_b(t)).\quad}
\]
Следовательно,
\[
 S(x,y,z)=(x,y,\sigma z+P(x,y))
\]
--- полиномиальный автоморфизм над $\mathbb Q$, переводящий
$\gamma_{a,b,c}$ в $\gamma_{a,b,c'}$; его обратный имеет вид
$(x,y,\sigma(z-P(x,y)))$.

\emph{Доказательство.}
Прообразные точки каждого узла можно записать в виде
\[
 s=2\cos\pi\left(\frac{k}{a}+\frac{\ell}{b}\right),\qquad
 t=2\cos\pi\left(\frac{k}{a}-\frac{\ell}{b}\right).
\]
По формуле разности косинусов
\[
 C_c(s)-C_c(t)
 =-4\sin\frac{ck\pi}{a}\sin\frac{c\ell\pi}{b}.
\]
При замене $c$ на $c'$ правая часть умножается на $\sigma$.
По лемме $C_{c'}-\sigma C_c\in A$. Все исходные коэффициенты
рациональны, поэтому алгоритм членения по ведущей степени даёт
$P\in\mathbb Q[x,y]$.\hfill$\square$

В частности,
\[
 H(a,b,c+2ab)\simeq_{\mathrm{alg},\mathbb Q}H(a,b,c).
\]
Это равенство алгебраических классов, а не только равенство типов
вещественных узлов.

\paragraph{Алгоритм нахождения $P$.}
В монической нормировке достаточно повторять следующую операцию.
Для старшего члена $\lambda t^d$ вычислить единственное
$s\in\{0,\ldots,a-1\}$, удовлетворяющее $sb\equiv d\pmod a$,
и положить $r=(d-sb)/a$. Если $r\geq0$, вычесть
$\lambda C_a^rC_b^s$ и добавить $\lambda x^ry^s$ к $P$.
Иначе степень $d$ --- пробел полугруппы; её член переносится в остаток.
Для полиномов из теоремы остаток обязан быть нулевым.
Это даёт одновременно алгоритм, тождество-сертификат и его прямую проверку.

\paragraph{Число представителей при фиксированной паре.}
Среди классов $c\pmod {2ab}$ ровно $2\varphi(a)\varphi(b)$ удовлетворяют
$\gcd(c,ab)=1$. Четыре независимых выбора знаков в теореме действуют
на этом множестве свободно: совпадение двух знаков по модулю $2a$
или $2b$ означало бы соответственно $a\mid c$ или $b\mid c$.
Поэтому число нужных представителей не превосходит
\[
 \boxed{\quad \frac{\varphi(a)\varphi(b)}2.\quad}
\]
Это верхняя оценка числа алгебраических классов. Некоторые представители
могут быть эквивалентны более общими автоморфизмами.

\subsection{Два полных бесконечных семейства как следствие базовых моделей}

Для $(a,b)=(3,4)$ достаточно двух представителей. В таблице
третья колонка означает остатки индекса $c$, а не значения полиномов.
\[
\begin{array}{c|c|l}
 (a,b)&\text{представитель}&c\pmod{2ab}\\\hline
 (3,4)&1&1,7,17,23\\
 (3,4)&5&5,11,13,19
\end{array}
\]
Таким образом, подтверждённое выпрямление $H(3,4,5)$ над полем $K$
немедленно даёт выпрямление \emph{всех} $H(3,4,c)$,
$\gcd(c,12)=1$, над тем же $K$.

Для $(a,b)=(3,5)$ достаточно четырёх представителей:
\[
\begin{array}{c|c|l}
 (a,b)&\text{представитель}&c\pmod{2ab}\\\hline
 (3,5)&1&1,11,19,29\\
 (3,5)&2&2,8,22,28\\
 (3,5)&4&4,14,16,26\\
 (3,5)&7&7,13,17,23
\end{array}
\]
Первые две строки выпрямляются над $\mathbb Q$ по уже известным результатам;
третья после перестановки координат есть $H(3,4,5)$;
четвёртая есть $H(3,5,7)$.
Поэтому подтверждённые выпрямления двух базовых моделей над $K$
дают выпрямление \emph{всех} $H(3,5,c)$,
$\gcd(c,15)=1$, над тем же $K$.

Например, следующие тождества проверяются точным раскрытием:
\[
\begin{aligned}
 C_7+C_1&=C_3C_4,\\
 C_{11}+C_5&=C_3(C_4^2-2),\\
 C_{11}+C_1&=(C_3^2-2)C_5,\\
 C_{14}+C_4&=(C_3^3-3C_3)C_5,\\
 C_{13}+C_7&=C_3(C_5^2-2).
\end{aligned}
\]
Здесь $C_j$ справа в скобках обозначает значение $C_j(t)$,
а не композицию, кроме явно выписанных обычных операций.

\paragraph{Вычислительные сертификаты.}
Файл \texttt{periodic\_shears.py} строит $P$ точной рациональной
арифметикой и непосредственно проверяет подстановкой 191 совместимое
сравнение индексов для пар
$(2,3),(3,4),(3,5),(4,5),(3,7),(5,7)$ при
$1\leq c\leq20$, $c<c'\leq40$.
Кроме того, проверены девять отдельных примеров, включая
$(a,b,c,c')=(7,9,11,137)$.
Результаты и явные полиномы сохранены в
\texttt{periodic\_shears\_results.json}.
Эти вычисления проверяют реализацию теоремы, но доказательство
теоремы не основано на конечном переборе.

\paragraph{Граница результата.}
Для всех гармонических вложений получены критерий вложения,
явное восстановление параметра и конечная редукция при фиксированной
паре индексов. Из этого не следует выпрямляемость произвольного
$H(a,b,c)$. В частности, тип узла
$H(5,6,7)$ совпадает с типом $H(4,5,7)$ по теореме 4.4
Koseleff--Pecker, но это топологическая изотопия и не даёт
полиномиального автоморфизма между комплексными вложениями.
Не следует также переносить утверждения о $H(3,5,7)$ на
$(t^3-3t,t^5-5t^3+4t,t^7-7t^5+14t^3-7t)$:
средняя координата этой последней модели отличается от $C_5$.

\subsection{Следствие о минимальной степени поля определения}

Пусть $\gamma$ определена над $\mathbb Q$ и её вещественный образ
задаёт нетривиальный длинный узел. Если полиномиальное выпрямление
$\gamma$ вместе с обратным отображением определено над числовым полем
$L$, то $L$ не имеет вещественных вложений.
В самом деле, применение любого вложения $L\hookrightarrow\mathbb R$
ко всем коэффициентам сохранило бы тождества обратимости и
$\Phi(\gamma(t))=(t,0,0)$, дав невозможное вещественное выпрямление.
Полиномиальный автоморфизм $\mathbb R^3$ является собственным
гомеоморфизмом и продолжается на одноточечную компактификацию,
поэтому не переводит нетривиальный узел в прямую.

Следовательно, $L$ вполне мнимо и $[L:\mathbb Q]$ чётно.
Если выпрямление построено над $\mathbb Q(\sqrt{-3})$, степень $2$
минимальна среди возможных степеней числовых полей определения.
Это утверждение не означает минимальности степеней компонент
полиномиального автоморфизма и не выделяет единственное допустимое
квадратичное поле.

\subsection{Фильтр исключительных слоёв}
\paragraph{Общий фильтр знаменателей для кубической первой координаты.}
Тот же аргумент применим к любой модификации
\[
 \beta=(x,y,m(x)z+F(x,y))
\]
на исходной модели с первыми двумя координатами $(C_3,C_b)$.
Если существует исправляющий автоморфизм $(x,H,W)$, у которого
$H(\beta\gamma(t))=t$ и каждый приведённый центр становится графиком
над $H$, то для каждого корня $r$ полинома $m$ необходимо
\[
 t=\lambda C_b(t)+\mu\quad\text{в }k[t]/(C_3(t)-r),\qquad\lambda\ne0.
\]
Этот вывод остаётся верен и для кратных корней $C_3(t)-r$:
он получается подстановкой в полиномиальное тождество по модулю
$C_3(t)-r$, а не только проверкой геометрических точек.

Положим $V_n(x)=U_n(x/2)$, $V_{-1}=0$, $V_0=1$,
$V_{n+1}=xV_n-V_{n-1}$. По индукции из чебышёвской рекурсии
получаем точные остатки
\[
\begin{aligned}
 C_{3n+1}(t)&\equiv -V_{n-1}(x)t^2+V_n(x)t+2V_{n-1}(x),\\
 C_{3n+2}(t)&\equiv V_n(x)t^2-V_{n-1}(x)t-2V_n(x)
 \pmod{C_3(t)-x}.
\end{aligned}
\]
Следовательно, приведённый многочлен $m_{\mathrm{red}}$ обязан делить
$V_{n-1}$ при $b=3n+1$ и $V_n$ при $b=3n+2$.
Соседние $V_n,V_{n-1}$ взаимно просты, так что линейный коэффициент
в остатке на разрешённых корнях не обнуляется.
Например, при $b=5$ и $b=7$ допустимы только корни $m$ в $x=0$;
при $b=8$ --- только $x=\pm1$; при $b=11$ --- только
$x=0,\pm\sqrt2$.
Это необходимое условие для указанного механизма переноса;
достаточность и условия на кратности корней здесь не установлены.

\section{Группы Кремоны и локально нильпотентные дифференцирования}

Ниже $k$ --- поле характеристики нуль, $A=k[x,y,z]$. Все утверждения в
этом разделе доказаны элементарно; мотивы, $K$-теория и гомологическая
алгебра не используются. Общие приёмы не объявляются новыми. В частности,
критерий подъёма параметра и устранение полюсов рациональных сдвигов
сопоставляются с работой Drew Lewis,
\emph{A unified approach to embeddings of a line in 3-space},
\url{https://arxiv.org/abs/2203.08178}: Lemma~1, Theorem~2 и Corollary~4.

\subsection{Точный критерий внутри класса переносов}

Пусть $(x,H,W)$ --- автоморфизм $\mathbb A^3_k$, сохраняющий $x$, и
\[
 d=H_yW_z-H_zW_y\in k^*,\qquad
 \beta=(x,p(x,y)+m(x)z,y),\qquad 0\ne m\in k[x].
\]
Обозначим
\[
 h=H\circ\beta,\quad w=W\circ\beta,\quad
 B=k[x]/(m),\quad G(y)=\overline{H(x,p(x,y),y)}\in B[y].
\]
Черта означает редукцию по $m$. В случае $m\in k^*$ перенос уже является
полиномиальным автоморфизмом; содержательный случай --- $\deg m>0$.

\paragraph{Утверждение.}
Дифференцирование
\[
 \delta_h=h_z\partial_y-h_y\partial_z
\]
всегда локально нильпотентно. Следующие четыре условия эквивалентны:
\begin{enumerate}
\item $G$ является координатой однопеременного кольца $B[y]$, то есть
подстановка $y\mapsto G(y)$ есть $B$-автоморфизм;
\item существуют $q\in A$ такие, что $(x,h,q)$ --- полиномиальный
автоморфизм;
\item у $\delta_h$ существует слайс: $\delta_h(s)=1$ для некоторого $s\in A$;
\item $\delta_h$ не имеет геометрических неподвижных точек, эквивалентно
$(h_y,h_z)=A$.
\end{enumerate}
В условии~1, если $G=\sum g_jy^j$, требуется и достаточно, чтобы $g_1$
был обратим в $B$, а все $g_j$ при $j\ge2$ были нильпотентны. Коэффициент
$g_0$ произволен. Поэтому проверка всех корней $m$ должна устанавливать
\emph{аффинность всей зависимости от $y$}, а не только ненулевую производную
в одной выбранной точке.

\paragraph{Доказательство локальной нильпотентности.}
В координатах $(x,H,W)$ дифференцирование
$\delta_H=H_z\partial_y-H_y\partial_z$ равно $-d\partial_W$.
После обращения $m$ подстановка $\beta^*$ является автоморфизмом и
\[
 \delta_h=-m\,\beta^*\delta_H(\beta^*)^{-1}.
\]
Множитель $m$ принадлежит ядру, поэтому правая часть локально
нильпотентна на $A_m$. Левая часть сохраняет $A$, а вложение $A\hookrightarrow
A_m$ инъективно; следовательно, её ограничение на $A$ также локально
нильпотентно. Это верно \emph{даже при невыполнении условия центра}.

\paragraph{Доказательство $1\Rightarrow2,3$.}
Пусть $f(T)\in B[T]$ --- обратная подстановка к $G$. Положим
\[
 J(y)=\overline{W(x,p(x,y),y)},\qquad \bar\varphi(T)=J(f(T)),
\]
и выберем произвольный полиномиальный подъём $\varphi\in k[x,T]$.
Тогда
\[
 q=\frac{w-\varphi(x,h)}{m}\in A.
\]
Обратное отображение строится без предположения об отсутствии кратных
корней. Пусть $(x,Y(x,T,V),Z(x,T,V))$ --- обратное к $(x,H,W)$.
При $V=mQ+\varphi(x,T)$ положим
\[
 y=Z(x,T,V),\qquad z=\frac{Y(x,T,V)-p(x,Z(x,T,V))}{m}.
\]
В кольце $B[T]$ имеем
\[
 (Y(x,T,\bar\varphi(T)),Z(x,T,\bar\varphi(T)))
   =(\bar p(x,f(T)),f(T)).
\]
Значит, последнее деление также точное. Обе композиции равны тождеству
после обращения $m$, следовательно, и в полиномиальном кольце. Кроме того,
\[
 \delta_h(w)=md,\qquad \delta_h(q)=d,
 \qquad \ker\delta_h=k[x,h].
\]
Таким образом, $q/d$ --- слайс. Именно необходимость проверки обратного
деления защищает от ошибочного перехода от бирациональности к автоморфизму.

\paragraph{Остальные импликации.}
Если $(x,h,q)$ --- автоморфизм, то его якобиан есть ненулевая константа,
так что $\delta_h(q)$ --- ненулевая константа: $2\Rightarrow3$.
Слайс несовместим с неподвижной точкой, откуда $3\Rightarrow4$.
Вне $m=0$ поле $\delta_h$ нигде не обращается в нуль, поскольку сопряжено
с ненулевой кратной постоянного поля. На геометрическом слое $x=\alpha$,
где $m(\alpha)=0$, оно равно
\[
 \delta_h\big|_{x=\alpha}=-G_\alpha'(y)\partial_z.
\]
Над алгебраическим замыканием полином $G_\alpha'$ не имеет корней тогда
и только тогда, когда он постоянен и отличен от нуля. В характеристике
нуль это равносильно тому, что $G_\alpha$ аффинен с ненулевым наклоном.
Для всех корней $m$ это ровно обратимость $g_1$ и нильпотентность высших
коэффициентов в $B$. Наконец, такая подстановка обратима: после обращения
линейной части остаток имеет коэффициенты в нильпотентном идеале, и
последовательное подставление заканчивается через конечное число шагов.
Получаем $4\Rightarrow1$.

\paragraph{Граница вывода.}
Это необходимый и достаточный критерий \emph{относительной} координаты
$(x,h)$ в указанном классе переносов. Он не является необходимым критерием
выпрямления произвольной кривой и здесь не утверждается как критерий того,
что $h$ есть абсолютная координата без сохранения $x$.
Если условие центра не выполнено, отвергаются именно этот перенос и слайс
именно этого $\delta_h$. Другие координаты и другие LND остаются возможны.
Без условия центра нельзя автоматически писать $\ker\delta_h=k[x,h]$:
например, $H=Y$, $p=xy$, $m=x$ дают $h=x(y+z)$ и
$\ker\delta_h=k[x,y+z]\supsetneq k[x,h]$.

\subsection{Конструктивная лемма о треугольных LND}

Рассмотрим $k[x,v,y]$ и
\[
 D=P(x)\partial_v+Q(x,v)\partial_y,\qquad 0\ne P\in k[x].
\]
Выберем $B_0(x,v)$ с $(B_0)_v=Q$ и положим
\[
 H=B_0(x,v)-P(x)y.
\]
Предположим, что в $k[x]/(P)$ постоянный по $v$ коэффициент $Q$ обратим,
а все остальные коэффициенты $Q$ нильпотентны. Тогда
$\bar B_0$ является автоморфизмом $(k[x]/(P))[v]$. Если $f(x,T)$ ---
полиномиальный подъём его обратного, то
\[
 W=\frac{v-f(x,H)}{P(x)}\in k[x,v,y],\qquad DH=0,\quad DW=1.
\]
Полиномиальность $W$ следует из $H\equiv B_0(v)\pmod P$. Обратные
формулы имеют вид
\[
 v=f(x,H)+P(x)W,\qquad
 y=\frac{B_0(x,f(x,H)+P(x)W)-H}{P(x)}.
\]
Второй числитель также делится на $P$, поскольку
$B_0(x,f(x,H))\equiv H\pmod P$. Поэтому $(x,H,W)$ --- автоморфизм,
\[
 k[x,v,y]=k[x,H,W],\qquad D=\partial_W,
 \qquad \ker D=k[x,H].
\]
Над алгебраически замкнутым полем гипотеза леммы эквивалентна отсутствию
неподвижных точек у этого треугольного LND: над каждым корнем $P$
полином $Q(\alpha,v)$ должен быть ненулевой константой.

В частном случае $Q=Q_1(x)v+R(x)$ критерий принимает простой вид
\[
 Q_1\in\sqrt{(P)},\qquad \gcd(P,R)=1.
\]
Для базового дифференцирования репозитория это объясняет роль
$P=(1-a)x^5$, $Q_1=2(1-a)x$, $R=-1+ax^2-ax^4$: старшая поправка
нильпотентна по модулю $x^5$, а линейная часть обратима. Это конструктивная
форма классического механизма, а не заявка на новую общую классификацию LND.

\paragraph{Небольшой полностью проверенный пример с кратным корнем.}
Для $m=x^2$, $p=y+xy^2$, $H=Y$, $W=Z$ центр имеет обратный
$f(T)=T-xT^2\pmod{x^2}$. Получаем
\[
 h=y+xy^2+x^2z,\qquad
 q=x^3z^2+2x^2y^2z+xy^4+2xyz+2y^3-z.
\]
Тогда $\delta_h=x^2\partial_y-(1+2xy)\partial_z$,
$\delta_h(h)=0$, $\delta_h(q)=1$, и обратные формулы к $(x,h,q)$ суть
\[
 y=T-xT^2+x^2Q,
\]
\[
 z=-Q^2x^3+2QT^2x^2-2QTx-Q-T^4x+2T^3.
\]
Обе композиции, равенства для LND и якобиан $-1$ проверены точно
скриптом \texttt{research/cremona\_lnd/verify\_example.py}.
Напротив, для $h=y^2+xz$ существует LND
$x\partial_y-2y\partial_z$, но она имеет неподвижную прямую $x=y=0$
и потому не имеет слайса. Это простой пример того, почему одного
построения LND недостаточно.

\subsection{Что даёт группа Кремоны и где возникают полюса}

Пусть $\gamma(t)=(f(t),g(t),h_0(t))$ --- замкнутое вложение. После общей
линейной замены двух координат можно добиться
$k(f,g)=k(t)$: это обычная теорема о примитивном элементе для конечного
сепарабельного расширения $k(t)/k(f)$. Поэтому найдётся
$r(x,y)\in k(x,y)$ с $r(f(t),g(t))=t$.
Положим
\[
 w=z-h_0(r(x,y))+r(x,y),\qquad
 F(x,y,z)=(x-f(w),y-g(w),w).
\]
Это бирациональное преобразование с явным обратным
\[
 x=X+f(s),\quad y=Y+g(s),\quad
 z=s-r(x,y)+h_0(r(x,y)).
\]
На плотном открытом подмножестве кривой $F\gamma(t)=(0,0,t)$.
Таким образом, бирациональное выпрямление здесь имеет элементарную
конструкцию. Оно ещё не даёт отображения, регулярного во всех точках
кривой, тем более глобального полиномиального автоморфизма.

Если числители и знаменатели построенного преобразования лежат в
$A_m$, задача полиномиальности сводится к точной делимости:
\[
 \frac{N}{m^e}\in A
 \quad\Longleftrightarrow\quad
 \operatorname{ord}_\pi(N)\ge e\operatorname{ord}_\pi(m)
 \quad\hbox{для всех неприводимых }\pi\mid m.
\]
При $m=\prod_\alpha(x-\alpha)^{r_\alpha}$ это равносильно занулению
всех $x$-производных числителя порядков $0,\ldots,er_\alpha-1$ при
$x=\alpha$. Проверка только значений в корнях недостаточна для кратных
корней. Требования применяются к прямому и обратному отображениям.

Предупреждение о LND существенно: рациональное сопряжение постоянного
поля может стать полиномиальным полем, не будучи локально нильпотентным.
Например, $\partial_u$ при $u=-1/x$ превращается в $x^2\partial_x$.
Безопасный вариант таков: если сопряжение является автоморфизмом $A_m$
над $k[x,m^{-1}]$, то сопряжённые $\partial_y,\partial_z$ локально
нильпотентны на $A_m$ и аннулируют $m$. Умножение на достаточно большую
степень $m$ делает их полиномиальными LND на $A$; два поля остаются
коммутирующими. Но возникающее вырождение на $m=0$ может уничтожить
слайсы. Поэтому очистка знаменателей LND сама по себе не решает задачу
выпрямления.

\subsection{Подъём параметра, стабилизация и действия на дополнении}

Для идеала $I=I(\gamma)$ замкнутость вложения означает наличие
$R\in A$ с $R\circ\gamma=t$. Выпрямляемость равносильна тому, что класс
$R+I$ содержит \emph{координату} $S$ кольца $A$. Действительно, если
$(S,U,V)$ --- координаты и $S\gamma=t$, то
\[
 (S,U-U(\gamma(S)),V-V(\gamma(S)))
\]
выпрямляет кривую; здесь $U(\gamma(t)),V(\gamma(t))$ сначала рассматриваются
как однопеременные полиномы, после чего подставляется $S$.
Обратное следует из первой компоненты выпрямляющего автоморфизма.
Это известный критерий; см. Lewis, Lemma~1.

После добавления \emph{одной} переменной любое такое вложение
выпрямляется явно. На $\mathbb A^4$ последовательно применяются
\[
 (x,y,z,u)\longmapsto(x,y,z,u+R(x,y,z))
\]
и
\[
 (x,y,z,u)\longmapsto(x-f(u),y-g(u),z-h_0(u),u).
\]
Они переводят $(\gamma(t),0)$ в $(0,0,0,t)$ и являются произведениями
элементарных полиномиальных автоморфизмов. Этот факт не позволяет убрать
добавленную переменную; именно такой переход и требует отдельного
доказательства в размерности три.

Если выпрямитель уже известен, обозначим его координаты $(S,U,V)$ так,
что $I=(U,V)$ и $S\gamma=t$. Сопряжение трёх частных производных даёт
коммутирующие LND $D_S,D_U,D_V$ с
\[
 D_i(S,U,V)=e_i.
\]
В частности, $D_S$ имеет слайс $S$, ядро $k[U,V]$ и орбиту $\gamma$.
Явное действие записывается через выпрямитель $\Phi$:
\[
 \lambda\cdot p=\Phi^{-1}(S(p)+\lambda,U(p),V(p)).
\]
Это даёт полноценное свободное алгебраическое действие $\mathbb G_a$,
а не только касательное поле вдоль кривой. Обратное утверждение о
\emph{произвольной} орбите произвольного LND здесь не используется.

Обычный инвариант Макар--Лиманова кольца регулярных функций дополнения
не различает рассматриваемые вложения. Для любой кривой
$C\subset\mathbb A^3$ коразмерности два
\[
 \Gamma(\mathbb A^3\setminus C,\mathcal O)=k[x,y,z].
\]
Элементарное объяснение: неприводимый множитель знаменателя рациональной
функции задаёт дивизор; такой дивизор не может целиком содержаться в $C$,
поэтому регулярность на дополнении исключает каждый такой множитель.
Следовательно, ML этого кольца всегда равен $k$.

Для сохранения информации о вложении можно рассматривать именно пару
$(A,I)$ и пересечение
\[
 \mathrm{ML}(A,I)=
 \bigcap_{\substack{D\in\operatorname{LND}(A)\\D(I)\subset I}}\ker D.
\]
Это определение фиксирует, что соответствующие действия сохраняют
удаляемую кривую. Для выпрямляемой кривой оно также равно $k$:
достаточно трёх сохраняющих $I$ LND
$D_S$, $U D_V$, $V D_U$, чьи ядра имеют пересечение $k$.
Таким образом, нетривиальность этого относительного пересечения была
бы препятствием, но его тривиальность не объявляется достаточной для
выпрямления. Вычисление всех сохраняющих $I$ LND --- отдельная трудная
задача; простая проверка конечного списка полей не устанавливает
нетривиальность полного пересечения.

\paragraph{Дополнительные первоисточники.}
K. Masuda, \emph{Homogeneous locally nilpotent derivations having slices
and embeddings of affine spaces}, J. Algebra 321 (2009), 1719--1733,
\url{https://doi.org/10.1016/j.jalgebra.2008.12.012}, устанавливает
связь вложений с определёнными однородными коммутирующими LND со
слайсами; важно сохранять условия на градуировку.
S. Kaliman, \emph{Free $\mathbb C_+$-actions on $\mathbb C^3$ are
translations}, \url{https://arxiv.org/abs/math/0207156}, даёт общий
результат о свободных действиях в размерности три. Его глубокая теорема
не нужна для доказанных выше специальных критериев.

\subsection{Элементарная структура базового дифференцирования}
Работаем над $K=\mathbb Q[a]/(a^2-a+1)$ и полагаем
\[
b=1-a,\qquad v=z+axy,\qquad
A(x)=-1+ax^2-ax^4,
\]
\[
q(x)=-2x^5+(6-2a)x^3-2bx,
\qquad h_0(x,v)=bxv^2+A(x)v+q(x).
\]
Прямое раскрытие скобок даёт более короткую запись базовой координаты:
\[
H(x,y,z)=h_0(x,z+axy)-bx^5y.
\]
Замена $(x,y,z)\leftrightarrow(x,v,y)$ полиномиальна в обе стороны.
Якобиево дифференцирование
\[
D=H_z\frac{\partial}{\partial y}-H_y\frac{\partial}{\partial z}
\]
в этих координатах имеет треугольный вид
\[
\boxed{D=bx^5\frac{\partial}{\partial v}
       +(2bxv+A(x))\frac{\partial}{\partial y}.}
\]
Здесь $D(x)=D(H)=0$, $D^2(v)=0$, $D^3(y)=D^3(z)=0$;
следовательно, $D$ локально нильпотентно. Это утверждение следует
не из численной проверки, а из указанного треугольного вида.

Пусть $f(x,T)$ --- полином из репозитория. Точные вычисления дают
\[
f(x,h_0(x,v))\equiv v\pmod{x^5},\qquad
h_0(x,f(x,T))\equiv T\pmod{x^5}.
\]
Эти тождества можно также получить конечным обращением
$h_0(x,v)\equiv-v\pmod{x}$ над $K[x]/(x^5)$.
Поэтому выражение
\[
W=\frac{a\bigl(v-f(x,H)\bigr)}{x^5}
\]
полиномиально. Так как $D(v)=bx^5$, $D(f(x,H))=0$ и $ab=1$, имеем
\[
\boxed{D(W)=1.}
\]
В частности, $W$ --- явный слайс. Для полноты координатности
предъявляется обратное преобразование:
\[
v=f(x,T)+bx^5V,\qquad
y=\frac{h_0(x,v)-T}{bx^5},\qquad z=v-axy.
\]
Последнее деление точно вследствие второго сравнения выше.
Эти формулы взаимно обратны к $(x,y,z)\mapsto(x,H,W)$ после
обращения $x$. Поскольку обе стороны полиномиальны, инъективность
локализации устанавливает обе композиции и на $x=0$.
Таким образом,
\[
K[x,y,z]=K[x,H,W],\qquad \ker D=K[x,H].
\]
Это конкретная реализация классического обращения по модулю
исключительного множителя; общая лемма не заявляется как новая.

Для гармонического переноса положим
\[
h_H=H(x,xz-y^2+2,y),\qquad
D_H=(h_H)_z\partial_y-(h_H)_y\partial_z.
\]
Если $q_H$ --- третий компонент относительного автоморфизма
репозитория, то его якобиан равен $+1$, поэтому
$D_H(q_H)=-1$. Следовательно, $-q_H$ есть слайс и
$\ker D_H=K[x,h_H]$. Независимое точное дифференцирование даёт
$D_H^3(y)=0$, $D_H^5(z)=0$, причём предыдущие итерации
$D_H^2(y)$ и $D_H^4(z)$ ненулевые. Это два конкретных
полиномиальных локально нильпотентных дифференцирования с
полностью явными координатами ядра и слайсами.

\section{Гипотеза ван ден Эссена и явные координаты ядер LND}

\subsection{Точная гипотеза и её логическое место}

Работа A.~van den Essen, P.~van Rossum, \emph{Triangular derivations related
to problems on affine $n$-space}, Proc. Amer. Math. Soc. \textbf{130}
(2002), 1311--1322, DOI
\href{https://doi.org/10.1090/S0002-9939-01-06212-8}{10.1090/S0002-9939-01-06212-8},
содержит соответствующую гипотезу как Conjecture~5.1, с.~1319.
Точная формулировка использует
\[
 f(u)=u^3-3u,\qquad g(u)=u^4-4u^2,\qquad h(u)=u^5-10u,
\]
\[
 A=\mathbb C[t,u,x,y,z],\qquad
 D=t\partial_u+f'(u)\partial_x+g'(u)\partial_y+h'(u)\partial_z.
\]
Утверждается, что $\ker D\not\cong_{\mathbb C}\mathbb C^{[4]}$.
Именно эту конкретную гипотезу рассматриваем здесь под названием
\emph{гипотеза E}. Множитель $t$ перед $\partial_u$ существенен.

В первоисточнике Theorem~4.1 и Corollary~4.2, с.~1317--1318, связывают
выпрямляемость вложения с полиномиальностью ядра. Наличие слайса
обсуждается в Theorem~3.1. Поэтому проверенный выпрямляющий автоморфизм
трилистника влечёт отрицание именно Conjecture~5.1. Это не является
доказательством общей проблемы сокращения: кандидат имеет размерность
$4$, а после добавления слайса получается полиномиальное кольцо
размерности $5$.

\subsection{Элементарный конструктивный перенос}

Следующее доказательство выписано непосредственно, чтобы зависимость от
проверяемого выпрямителя была полностью прозрачной. Сам принцип переноса
не заявляется как новый: он согласуется с указанным Theorem~4.1.

Пусть $k$ --- поле характеристики нуль,
$\gamma(u)=(f_1(u),f_2(u),f_3(u))$, и дан полиномиальный автоморфизм
\[
 F=(P,Q,R),\qquad F(\gamma(u))=(u,0,0),\qquad G=F^{-1}.
\]
Пусть $0\ne b(t)\in k[t]$. Рассмотрим
\[
 D_b=b(t)\partial_u+\sum_{i=1}^3f_i'(u)\partial_{x_i}
 \quad\text{на }A=k[t,u,x_1,x_2,x_3].
\]
Обозначим $V_i=f_i(u)-b(t)x_i$, $V=(V_1,V_2,V_3)$ и положим
\begin{equation}\label{eq:lnd-coordinates}
 p=P(V),\qquad q=\frac{Q(V)}{b(t)},\qquad
 r=\frac{R(V)}{b(t)},\qquad s=\frac{u-P(V)}{b(t)}.
\end{equation}
Тогда все четыре выражения являются полиномами и
\begin{equation}\label{eq:lnd-conjugacy}
 A=k[t,p,q,r,s],\qquad D_b=\partial_s,\qquad
 \ker D_b=k[t,p,q,r].
\end{equation}

\paragraph{Доказательство полиномиальности.}
Для любого полинома $L$ разность $L(\gamma(u)-b(t)X)-L(\gamma(u))$
делится на $b(t)$. Применив это к $P,Q,R$ и использовав
$P(\gamma)=u$, $Q(\gamma)=R(\gamma)=0$, получаем точную делимость
числителей в \eqref{eq:lnd-coordinates}. Это рассуждение происходит
в кольце полиномов; простота корней $b$ не требуется.

\paragraph{Явная обратная замена.}
Считая $t,p,q,r,s$ независимыми переменными, определим
\begin{equation}\label{eq:lnd-inverse}
 u=p+b(t)s,\qquad
 x_i=\frac{f_i(p+b(t)s)-G_i(p,b(t)q,b(t)r)}{b(t)}.
\end{equation}
Из $G(p,0,0)=\gamma(p)$ следует, что числитель также делится на
$b(t)$. В частности, формулы действуют при кратных корнях $b$.
Подстановка \eqref{eq:lnd-inverse} даёт
$V=G(p,b(t)q,b(t)r)$; последующее применение $F$ восстанавливает
$(p,b(t)q,b(t)r)$. Обратная композиция столь же непосредственно
восстанавливает $u$ и $X$. Следовательно, это взаимно обратные
полиномиальные замены координат.

Наконец, $D_b(t)=D_b(V_i)=0$, поэтому $D_b(p)=D_b(q)=D_b(r)=0$,
а $D_b(s)=1$. Это доказывает \eqref{eq:lnd-conjugacy}.

\paragraph{Что именно даёт применение.}
Для каждой модели с удостоверенным автоморфизмом $F$ получаем не только
абстрактную полиномиальность ядра, но явное сопряжение бесконечного
семейства $D_b$ с переносом. Ранг каждого такого ненулевого LND равен
$1$: в одной системе координат он уничтожает четыре координаты и
дифференцирует пятую. Новое применение к конкретной модели следует
признавать только после проверки обоих полиномиальных автоморфизмов
$F$ и $G$; проверка одного значения $F(\gamma(u))$ недостаточна.

\subsection{Слайс существует и без выпрямления}

Для трилистника имеется простой полином восстановления параметра
\[
 F_0(x,y,z)=yz-x^3-5xy+2z-7x,\qquad
 F_0(f(u),g(u),h(u))=u.
\]
Отсюда
\[
 s_0=\frac{u-F_0(f(u)-tx,g(u)-ty,h(u)-tz)}{t}
\]
является полиномом и $D(s_0)=1$. В компактной форме
\begin{align*}
 s_0={}&(-3f(u)^2-5g(u)-7)x+(h(u)-5f(u))y+(g(u)+2)z\\
       &+t\bigl(3f(u)x^2+5xy-yz\bigr)-t^2x^3.
\end{align*}
Эти тождества проверяются точной арифметикой в прилагаемом скрипте.

Однако $F_0$ не является координатой аффинного трёхмерного пространства:
\[
 \nabla F_0=(-3x^2-5y-7,\ z-5x,\ y+2)
\]
обращается в нуль в точках $(1,-2,5)$ и $(-1,-2,-5)$.
Координата полиномиального автоморфизма не имеет критических точек,
поскольку соответствующая строка обратимой матрицы Якоби не может
быть нулевой. Таким образом, полиномиальное восстановление параметра
и наличие слайса сами по себе не заменяют построения $F$.

Для нормированных полиномов Чебышёва $C_n(u)=2T_n(u/2)$ имеем
\[
 (C_3,C_4,C_5)=(f,g+2,h-5f).
\]
В этих аффинно изменённых координатах полином восстановления принимает
особенно простой вид
\[
 P_0(x,y,z)=yz-x^3+3x,\qquad P_0(C_3(u),C_4(u),C_5(u))=u.
\]
Это также непосредственно следует из
$C_4C_5=C_9+C_1$ и $C_3(C_3(u))=C_9(u)$.
Полином $P_0$ имеет критические точки $(1,0,0)$ и $(-1,0,0)$;
следовательно, даже столь короткая формула восстановления параметра
не является первой координатой выпрямляющего автоморфизма.

\subsection{Ещё два конкретных следствия и линеаризация действия}

Пусть $F=(P,Q,R_3)$ --- удостоверенный выпрямитель трилистника над
$K=\mathbb Q(\sqrt{-3})$, а $b(t)=t$. Приведённые формулы дают
$B=\ker D=K[t,p,q,r]$ и $A=B[s]$. Они непосредственно устраняют также
два связанных кандидата из того же первоисточника: Conjectures~5.2 и~5.3,
с.~1320. Ниже не используются методы гомологической алгебры.

\paragraph{Гиперповерхность из Conjecture~5.2.}
Обозначим через $\sigma=s_0$ исходный слайс, построенный из $F_0$.
Поскольку $D(\sigma)=D(s)=1$, разность
\[
 h_0=\sigma-s=\frac{P(V)-F_0(V)}t
\]
принадлежит $B$. Поэтому замена $s\mapsto s+h_0$ является
треугольным автоморфизмом $B[s]$, и
\[
 A=K[t,p,q,r,\sigma],\qquad A/(\sigma)=K[t,p,q,r]\cong K^{[4]}.
\]
Таким образом, $\sigma$ --- координата пятимерного пространства.
Это совместимо с некоординатностью $F_0$ на трёхмерном пространстве:
речь идёт о разных полиномах в разных кольцах.

\paragraph{Расширенное кольцо Риса из Conjecture~5.3.}
Положим $R_0=K[x,y,z]=K[P,Q,R_3]$ и $I=I_\gamma=(Q,R_3)$.
Рассматриваемое кольцо, определённое как подкольцо $R_0[t,t^{-1}]$,
имеет явное описание
\[
 R_0[t,I/t]=K[t,P,Q/t,R_3/t]\cong K^{[4]}.
\]
Действительно, старые координаты $Q,R_3$ восстанавливаются умножением
новых на $t$, а алгебраическая независимость новых четырёх образующих
следует после обращения $t$. После расширения скаляров до $\mathbb C$
оба изоморфизма сохраняются, давая отрицание именно этих двух
конкретных гипотез, как и Conjecture~5.1.

\paragraph{Естественное действие мультипликативной группы.}
Зададим действие подстановками
\[
 t\longmapsto\lambda t,\qquad
 u\longmapsto u,\qquad X\longmapsto\lambda^{-1}X
 \quad(\lambda\in K^*).
\]
Вектор $V=\gamma(u)-tX$ инвариантен. Поэтому на ядре $B$ новые
координаты $(t,p,q,r)$ имеют веса $(1,0,-1,-1)$. Следовательно,
это конкретное действие $\mathbb G_m$ на $\operatorname{Spec}B$
линеаризуется явной заменой координат. Общая проблема линеаризации
действий на аффинных пространствах отсюда не решается.

\subsection{Дихотомия Макар--Лиманова для связанного кольца}

Есть элементарное следствие, пригодное для дальнейшего поиска.
Пусть $B=\ker D_\gamma$ для
\[
 D_\gamma=t\partial_u+\sum_{i=1}^n f_i'(u)\partial_{x_i}
 \quad\text{на }k[t,u,x_1,\ldots,x_n].
\]
Тогда
\begin{equation}\label{eq:ml-dichotomy}
 \operatorname{ML}(B)\in\{k,\ k[t]\}.
\end{equation}
Здесь $\operatorname{ML}(B)$ --- пересечение ядер всех локально
нильпотентных $k$-дифференцирований $B$.

Действительно, $\partial_{x_i}$ коммутируют с $D_\gamma$, а потому
ограничиваются до LND на $B$. Следовательно,
\[
 \operatorname{ML}(B)\subseteq
 \bigcap_i\ker(\partial_{x_i}|_B)
 =B\cap k[t,u]=k[t].
\]
Если $\operatorname{ML}(B)$ содержит непостоянный полином $a(t)$, то
для всякого LND $\delta$ на $B$ имеем
$0=\delta(a(t))=a'(t)\delta(t)$. Кольцо является областью целостности,
характеристика равна нулю, поэтому $\delta(t)=0$. Значит $k[t]$
содержится в $\operatorname{ML}(B)$, что доказывает дихотомию.

Она задаёт два конкретных исследовательских направления. Достаточно
найти один LND $\delta$ на $B$ с $\delta(t)\ne0$, чтобы получить
$\operatorname{ML}(B)=k$. Если удастся доказать, что все LND уничтожают
$t$, то $\operatorname{ML}(B)=k[t]$, кольцо $B$ не полиномиально, и
соответствующее вложение невыпрямляемо. Обратная импликация
\emph{тривиальный ML $\Rightarrow$ полиномиальное кольцо} здесь не
утверждается. При известном выпрямителе формулы выше дают нужный LND
$\partial_t$ в координатах $(t,p,q,r)$.

\subsection{Сопряжённые LND на исходном пространстве}

Для проверенного $F=(P,Q,R)$ перенос координатных дифференцирований
даёт три коммутирующих LND $\delta_P,\delta_Q,\delta_R$ на
$k[x,y,z]$. В частности,
\[
 \delta_P(P)=1,\quad \delta_P(Q)=\delta_P(R)=0,
 \qquad\delta_P|_{\gamma}=\frac{d}{du}.
\]
Если $d=\det JF\in k^*$, то
\[
 \delta_P(a)=d^{-1}\det\frac{\partial(Q,R,a)}{\partial(x,y,z)}.
\]
Идеал кривой имеет вид $I_\gamma=(Q,R)$. Дифференцирования
$\delta_P$, $Q\delta_R$ и $R\delta_Q$ сохраняют этот идеал, и
пересечение их ядер равно $k$. Поэтому обычный инвариант
Макар--Лиманова самой $\mathbb A^1$ или окружающей $\mathbb A^3$
не различает рассматриваемые вложения. Содержательны условия,
связывающие LND с идеалом вложенной кривой, либо инварианты
связанного кольца $B$.

\subsection{Ограничения интерпретации и источники}

Невыпрямляемость вещественного узла не противоречит существованию
комплексного автоморфизма. Из комплексной полиномиальности ядра
нельзя без отдельного аргумента выводить вещественную или рациональную
полиномиальность. Также нельзя обращать импликацию
\emph{выпрямляемость $\Rightarrow$ полиномиальность ядра}, игнорируя
дополнительные условия на градуировку.

У K.~Masuda, \emph{Equivariant embeddings and semi-invariant locally
nilpotent derivations}, Theorems~1.1--1.2, рассматривается именно
эквивариантное соответствие и полиномиальные координаты с заданными
весами; это полезно для дальнейшего исследования симметрий гармонических
моделей. Доступный авторский текст:
\url{https://sci-tech.ksc.kwansei.ac.jp/~kayo/cmn/papers/eembed.pdf}.

Препринт S.~Oda \url{https://arxiv.org/abs/1201.4198}, появляющийся
в поиске как прежнее отрицательное решение гипотезы E, отозван автором
9 мая 2020 года с указанием ошибочности. Его утверждение не используется.

Первичный текст van den Essen--van Rossum доступен по ссылке издателя
\url{https://www.ams.org/journals/proc/2002-130-05/S0002-9939-01-06212-8/S0002-9939-01-06212-8.pdf};
проверенная копия того же журнального PDF:
\url{https://scispace.com/pdf/triangular-derivations-related-to-problems-on-affine-n-space-3val4k6v2j.pdf}.

\subsection{Минимальная степень слайса: точный результат, независимый от выпрямления}

Рассмотрим над произвольным полем $k$ характеристики нуль дифференцирование
\[
 D=t\partial_u+(3u^2-3)\partial_x+(4u^3-8u)\partial_y
 +(5u^4-10)\partial_z
 \quad\text{на }k[t,u,x,y,z].
\]
Слайсом называется полином $s$ такой, что $D(s)=1$.
Полная степень ниже является обычной полной степенью по всем пяти переменным
$t,u,x,y,z$, включая $t$.

\paragraph{Теорема.}
Минимальная полная степень полиномиального слайса для $D$ равна $7$.
Утверждение не использует выпрямитель из репозитория. Здесь не заявляется,
что этот результат отсутствует в литературе: установлена его самостоятельная
точная проверка с проверяемым сертификатом нижней границы.

\paragraph{Верхняя граница.}
Положим $f=u^3-3u$, $g=u^4-4u^2$, $h=u^5-10u$. Тогда
\begin{align*}
 \sigma={}&(g+2)z+(h-5f)y-(3f^2+5g+7)x\\
 &+t(3fx^2+5xy-yz)-t^2x^3.
\end{align*}
Иными словами,
\begin{align*}
 \sigma={}&-(3u^6-13u^4+7u^2+7)x
 +(u^5-5u^3+5u)y+(u^4-4u^2+2)z\\
 &+t\bigl(3(u^3-3u)x^2+5xy-yz\bigr)-t^2x^3.
\end{align*}
У этого полинома $15$ ненулевых мономов, полная степень $7$,
и непосредственное дифференцирование даёт $D(\sigma)=1$.
Полезное происхождение формулы: если
\[
 F_0=yz-x^3-5xy+2z-7x,\qquad
 V=(f,g,h)-t(x,y,z),
\]
то $F_0(f(u),g(u),h(u))=u$ и $\sigma=(u-F_0(V))/t$.
Последнее выражение является точным полиномиальным частным, а не
рациональной функцией с неустранённым полюсом.

\paragraph{Элементарное доказательство нижней границы.}
Введём вес: $\operatorname{wt}(t)=1$, $\operatorname{wt}(u)=0$,
$\operatorname{wt}(x)=\operatorname{wt}(y)=\operatorname{wt}(z)=-1$.
Дифференцирование $D$ повышает вес на единицу. Поэтому если бы существовал
слайс полной степени не выше $6$, его компонента веса $-1$ также была бы
слайсом той же степени. Обозначим $X=(x,y,z)$ и
$D_0=f'(u)\partial_x+g'(u)\partial_y+h'(u)\partial_z$.
Компонента веса $-1$ имеет вид
\[
 s=s_0(u,X)+t s_1(u,X)+t^2s_2(u,X),
\]
где $s_i$ однороден по $X$ степени $i+1$, а
$\deg_u s_i\leq 5-2i$. Сравнение коэффициентов при степенях $t$ даёт
\[
 D_0s_0=1,\qquad \partial_us_0+D_0s_1=0,\qquad
 \partial_us_1+D_0s_2=0,\qquad \partial_us_2=0.
\]
Последнее равенство означает, что $s_2$ --- кубическая форма от $X$
с постоянными коэффициентами. В третьем равенстве степень
$\partial_us_1$ по $u$ не выше $2$. Коэффициенты при $u^4$ и $u^3$
поэтому дают $\partial_zs_2=\partial_ys_2=0$, то есть $s_2=cx^3$.
Интегрируя третье равенство, получаем
\[
 s_1=-3c(u^3-3u)x^2+K_2(X),
\]
где $K_2$ --- постоянная квадратичная форма. Коэффициент при $u^5$
во втором равенстве равен $-18cx$, поскольку $\deg_u\partial_us_0\leq4$
и $\deg_u D_0K_2\leq4$. Значит, $c=0$. Следовательно,
\[
 s_1=K_2(X),\qquad
 s_0=L_1(X)-\gamma(u)\cdot\nabla K_2(X)
\]
для некоторой линейной формы $L_1$. Первое равенство теперь принимает вид
\[
 \frac{d}{du}\left(L_1(\gamma(u))-K_2(\gamma(u))\right)=1.
\]
Так как $\gamma(0)=0$, получаем
$L_1(\gamma(u))-K_2(\gamma(u))=u$.
Но нечётная часть ограничения любого полинома степени не выше $2$
на $\gamma$ лежит в линейной оболочке $f,h,fg,gh$.
Их степени $3,5,7,9$ попарно различны, поэтому никакая их линейная
комбинация не может равняться $u$. Противоречие доказывает нижнюю границу.
Доказательство работает над любым полем характеристики нуль.

\paragraph{Независимый вычислительный сертификат нижней границы.}
Для монома $m$ обозначим через $[m]A$ его коэффициент в полиноме $A$.
Определим $\mathbb Q$-линейный функционал
\[
 L(A)=\sum_{m\in\mathcal S} c_m[m]A,
\]
где все $31$ ненулевой вес $c_m$ приведены в таблице ниже;
веса остальных мономов равны нулю.
Точная проверка по правилу Лейбница устанавливает
\[
 L(1)=1,\qquad L(Dm)=0
 \quad\text{для каждого монома }m\text{ полной степени }\leq 6.
\]
Таких мономов ровно $\binom{11}{5}=462$.
Для любого полинома $s$ степени не выше $6$ отсюда следует $L(Ds)=0$.
Следовательно, равенство $Ds=1$ невозможно.
Это вместе с явным $\sigma$ доказывает теорему.

\begin{center}
\renewcommand{\arraystretch}{1.25}
\begin{tabular}{r r | r r}
$m$ & $c_m$ & $m$ & $c_m$ \\
\hline
$1$ & $1$ & $t^{2}z^{2}$ & $\frac{13375}{126}$ \\
$u^{2}$ & $1$ & $t^{2}y^{2}$ & $\frac{128}{63}$ \\
$u^{4}$ & $2$ & $t^{2}xz$ & $\frac{725}{42}$ \\
$u^{6}$ & $\frac{34}{7}$ & $t^{2}x^{2}$ & $\frac{27}{7}$ \\
$u^{8}$ & $\frac{116}{9}$ & $t^{2}uyz$ & $\frac{50}{9}$ \\
$tuz$ & $-\frac{50}{7}$ & $t^{2}uxy$ & $\frac{5}{3}$ \\
$tux$ & $-\frac{3}{2}$ & $t^{2}u^{2}z^{2}$ & $\frac{3875}{42}$ \\
$tu^{2}y$ & $-\frac{8}{7}$ & $t^{2}u^{2}y^{2}$ & $-\frac{152}{21}$ \\
$tu^{3}z$ & $-\frac{100}{9}$ & $t^{2}u^{2}xz$ & $\frac{995}{42}$ \\
$tu^{3}x$ & $-\frac{15}{7}$ & $t^{2}u^{2}x^{2}$ & $\frac{27}{7}$ \\
$tu^{4}y$ & $-\frac{160}{63}$ & $t^{2}u^{3}yz$ & $\frac{1430}{63}$ \\
$tu^{5}z$ & $-\frac{6275}{252}$ & $t^{2}u^{3}xy$ & $\frac{143}{21}$ \\
$tu^{5}x$ & $-\frac{57}{14}$ & $t^{2}u^{4}z^{2}$ & $\frac{13375}{63}$ \\
$tu^{6}y$ & $-\frac{149}{63}$ & $t^{2}u^{4}y^{2}$ & $-\frac{304}{21}$ \\
$tu^{7}z$ & $-\frac{12575}{252}$ & $t^{2}u^{4}xz$ & $\frac{1625}{42}$ \\
$tu^{7}x$ & $-\frac{859}{84}$ &  &  \\
\end{tabular}
\end{center}

Все веса рациональны. Продолжение $L$ по $k$-линейности доказывает ту же
нижнюю границу над любым полем $k$ характеристики нуль, в частности
над $\mathbb Q(\sqrt{-3})$ и над $\mathbb C$.
Нижняя граница относится ко всем полиномиальным слайсам, а не только
к выбранному анзацу или к полиномам, полученным из $F_0$.

\paragraph{Воспроизводимость.}
Файл \texttt{verify\_minimal\_slice.py} использует только стандартную
библиотеку Python и точные рациональные числа. Он читает
\texttt{no\_slice\_degree6\_certificate.json}, проверяет все $462$ равенства,
нормировку $L(1)=1$ и непосредственно вычисляет $D(\sigma)$.
Воспроизводимый сертификат не зависит от решения системы линейных
уравнений, от SymPy и от формул выпрямления. Дополнительно выполнена
независимая проверка с другим кодом правила Лейбница; объединение
мономов в образах $Dm$ имеет $812$ элементов и максимальную степень $9$.

\paragraph{Связь с явными координатами ядра.}
Если уже установлен выпрямитель $F=(P,Q,R)$ и построены
$p=P(V)$, $q=Q(V)/t$, $r=R(V)/t$, $s=(u-p)/t$, то компактный слайс
$\sigma$ можно использовать вместо $s$:
\[
 \sigma-s=\frac{p-F_0(F^{-1}(p,tq,tr))}{t}\in\ker D.
\]
Это полиномиальное выражение в $t,p,q,r$, поэтому
$(t,p,q,r,\sigma)$ также является системой полиномиальных координат.
В отличие от утверждения о минимальной степени, этот последний вывод
использует проверенный выпрямитель.

\subsection{Гармоническая восьмёрка: минимальная степень слайса равна одиннадцати}

Положим
\[
 f=C_3=u^3-3u,\quad g=C_5=u^5-5u^3+5u,\quad
 h=C_7=u^7-7u^5+14u^3-7u,
\]
и рассмотрим над полем $k$ характеристики нуль
\[
 D_H=t\partial_u+f'(u)\partial_x+g'(u)\partial_y+h'(u)\partial_z.
\]
Здесь полная степень снова считается по всем пяти переменным $t,u,x,y,z$.

\paragraph{Теорема.}
Минимальная полная степень полиномиального слайса для $D_H$ равна $11$.
Доказательство ниже элементарно и не использует выпрямитель из репозитория.
Литературная новизна этого конкретного результата не заявляется.

\paragraph{Кубическая ретракция параметра.}
Тождества Чебышёва
\[
 C_3C_7-C_5^2=C_4-2,\qquad C_3C_4=C_7+u
\]
дают
\[
 R_H(x,y,z)=x^2z-xy^2+2x-z,\qquad R_H(f(u),g(u),h(u))=u.
\]
Степень $3$ этой ретракции минимальна: поскольку $f,g,h$ нечётны,
нечётная часть ограничения любого полинома степени не выше $2$
лежит в линейной оболочке $f,g,h$. Их различные степени $3,5,7$
не позволяют получить $u$. При этом $R_H$ не является координатой
объемлющего пространства: его градиент обращается в нуль в точках
$(1,0,-1)$ и $(-1,0,1)$.

\paragraph{Явный слайс степени одиннадцать.}
Положим $V=(f,g,h)-t(x,y,z)$. Точное полиномиальное частное
$\sigma_H=(u-R_H(V))/t$ имеет компактную форму
\begin{align*}
 \sigma_H={}&(2fh-g^2+2)x-2fg\,y+(f^2-1)z\\
 &+t\left(fy^2-2fxz+2gxy-hx^2\right)
 +t^2\left(x^2z-xy^2\right).
\end{align*}
После раскрытия скобок это $27$ ненулевых мономов полной степени $11$;
коэффициент при $u^{10}x$ равен $1$.
Из $D_H(V)=0$ сразу следует $D_H(\sigma_H)=1$;
тождество также проверено прямым точным дифференцированием.

\paragraph{Доказательство нижней границы.}
Используем вес $\operatorname{wt}(t)=1$, $\operatorname{wt}(u)=0$,
$\operatorname{wt}(X)=-1$, где $X=(x,y,z)$.
Дифференцирование повышает вес на единицу, поэтому гипотетический
слайс полной степени не выше $10$ можно заменить его компонентой веса $-1$:
\[
 s=\sum_{i=0}^{4}t^i s_i(u,X),\qquad
 \deg_Xs_i=i+1,\quad \deg_us_i\leq9-2i,
\]
причём каждое $s_i$ однородно по $X$.
Обозначим $D_0=f'\partial_x+g'\partial_y+h'\partial_z$.
Сравнивая коэффициенты в $D_Hs=1$, получаем
\[
 D_0s_0=1,\qquad \partial_us_i+D_0s_{i+1}=0\ (0\leq i<4),
 \qquad \partial_us_4=0.
\]

Последнее равенство делает $s_4$ постоянной формой степени $5$ от $X$.
В уравнении $\partial_us_3+D_0s_4=0$ первая часть имеет степень по $u$
не выше $2$. Сравнение коэффициентов при $u^6$ и $u^4$ даёт
$\partial_zs_4=\partial_ys_4=0$, откуда $s_4=cx^5$ и
\[
 s_3=-5cf x^4+K_4(X).
\]
В следующем уравнении $\deg_u\partial_us_2\leq4$.
Коэффициент при $u^6$ даёт $\partial_zK_4=0$;
коэффициент при $u^5$ равен $-60cx^3$, поскольку $D_0K_4$
содержит только чётные степени $u$. Значит, $c=0$ и $K_4=K_4(x,y)$.
Интегрирование теперь даёт
\[
 s_2=-f\,\partial_xK_4-g\,\partial_yK_4+K_3(X).
\]
В уравнении $\partial_us_1+D_0s_2=0$ первая часть имеет степень по $u$
не выше $6$. Коэффициенты при $u^9$ и $u^7$ последовательно дают
$\partial_{yy}K_4=\partial_{xy}K_4=0$. Так как $K_4$ однороден степени $4$,
он имеет вид $K_4=a x^4$. Следовательно,
\[
 s_2=-4af x^3+K_3(X),\qquad
 s_1=6af^2x^2-\gamma(u)\cdot\nabla K_3(X)+K_2(X).
\]
В уравнении $\partial_us_0+D_0s_1=0$ первая часть имеет степень по $u$
не выше $8$. Члены при $u^{13}$, $u^{11}$ и $u^9$ дают соответственно
\[
 \partial_{zz}K_3=0,\quad \partial_{yz}K_3=0,\quad
 \partial_{yy}K_3=-2\partial_{xz}K_3.
\]
Таким образом,
\[
 K_3=d x^3+e x^2y+k(x^2z-2xy^2).
\]
Ещё одно интегрирование даёт
\[
 s_0=-4af^3x+\tfrac12\gamma(u)^{\mathsf T}
      \operatorname{Hess}K_3(X)\gamma(u)
      -\gamma(u)\cdot\nabla K_2(X)+K_1(X),
\]
где $K_j$ --- однородная форма степени $j$ с постоянными коэффициентами.
По симметрии коэффициентов кубической формы равенство $D_0s_0=1$
теперь эквивалентно
\[
 \frac{d}{du}\left(K_1(\gamma)-K_2(\gamma)+K_3(\gamma)-af^4\right)=1.
\]
Так как $\gamma(0)=0$, выражение в скобках должно равняться $u$.
Его нечётная часть лежит в линейной оболочке
\[
 f,\quad g,\quad h,\quad f^3,\quad f^2g,\quad f^2h-2fg^2.
\]
Степени этих шести полиномов равны соответственно $3,5,7,9,11,13$;
старший коэффициент последнего равен $-1$.
Последовательное сравнение старших коэффициентов исключает возможность
получить из их линейной комбинации полином $u$. Противоречие доказывает
нижнюю границу и теорему.

\paragraph{Независимый точный сертификат.}
Помимо аналитического доказательства, построен рациональный функционал
извлечения коэффициентов со $117$ ненулевыми весами, для которого
$L(1)=1$ и $L(D_Hm)=0$ на всех $\binom{15}{5}=3003$ мономах полной
степени не выше $10$. Все веса находятся в файле
\texttt{no\_harmonic\_slice\_degree10\_certificate.json}.
Автономный \texttt{verify\_harmonic\_minimal\_slice.py} проверяет
каждое из этих равенств с помощью стандартной библиотеки Python и
точных рациональных чисел; он также проверяет компактную формулу
$\sigma_H$ и все $27$ её коэффициентов.
Независимый второй проверяющий код дал тот же результат.
Объединение мономов всех образов $D_Hm$ имеет $6017$ элементов
и максимальную степень $15$.

\section{Точные вычисления для гармонических моделей}
Положим $C_n(t)=2T_n(t/2)$. Все приведённые ниже проверки выполнены
над точными рациональными коэффициентами в SymPy 1.14.0. Непрохождение
ограниченного алгоритма не означает невыпрямляемость.

\subsection{Перебор понижающих элементарных преобразований}
Если $c=ra+sb$ при $r,s\geq1$, то
\[
 (x,y,z)\longmapsto(x,y,C_r(x)C_s(y)-z)
\]
есть полиномиальная инволюция и переводит $(C_a,C_b,C_c)$ в
$(C_a,C_b,C_{|ra-sb|})$. Поскольку $0<|ra-sb|<c$, сумма индексов
строго убывает. Скрипт \texttt{reduction\_scan.py} перебирает \emph{все}
представления $c=ra+sb$ на каждом шаге, сортирует новую тройку и
мемоизирует множество всех достижимых терминальных троек. Поэтому
это полный перебор строго понижающих путей указанного вида, а не
жадный выбор единственного представления.

Тройка с индексом 1 непосредственно выпрямляется над $\mathbb Q$.
Достижение $(3,4,5)$ или $(3,5,7)$ даёт выпрямление над
$K=\mathbb Q(\sqrt{-3})$ при использовании проверенных сертификатов
репозитория. Других преобразований алгоритм не использует.
\begin{center}
\begin{tabular}{r|r|r|r|r}
Максимум & Всего & Достигнут индекс 1 & Базы 345/357 & Не покрыто \\
\hline
20 & 297 & 82 & 13 & 202\\
50 & 5382 & 734 & 72 & 4576\\
100 & 44142 & 3209 & 237 & 40696
\end{tabular}
\end{center}
Здесь перечислены строго возрастающие попарно взаимно простые тройки
$2\leq a<b<c\leq N$. Категории не пересекаются: при достижении
индекса 1 тройка учитывается в первом положительном столбце.
Это нижняя граница числа выпрямляемых моделей, а не классификация
остальных моделей. В частности, алгоритм не исследует пути с временным
увеличением степеней и произвольные полиномиальные автоморфизмы.
Для $a=3$ соответствующие числа при $N=100$ равны
$1409=425+142+842$. Первые непокрытые ядра при $a=3$:
$(3,7,8),(3,7,11),(3,8,13),(3,10,11),(3,10,17)$.
Другой важный первый пример: $(4,5,7)$.

\subsection{Матрица при кубической проекции}
Пусть $x=C_3(t)$, $s=C_2(t)=t^2-2$ и
$U_j=U_j(x/2)$ --- многочлен Чебышёва второго рода, причём $U_{-1}=0$.
Из рекурсии $C_{n+3}=xC_n-C_{n-3}$ следует
\[
 C_{3q+1}=U_qt-U_{q-1}s,\qquad
 C_{3q+2}=-U_{q-1}t+U_qs.
\]
Таким образом, $C_b=A_b(x)t+B_b(x)s$, и любые две координаты
дают явную матрицу $N_{b,c}$ над $\mathbb Q[x]$.
Для первых моделей её определители равны
\[
\begin{array}{c|c}
(b,c)&\det N_{b,c}\\\hline
(4,5)&x^2-1\\
(5,7)&-x(x^2-2)\\
(7,8)&x^4-3x^2+1\\
(7,11)&x(x^2-1)(x^2-3)
\end{array}
\]
Поэтому рациональное восстановление $(t,s)$ элементарно. Обращение
определителя не доказывает существование полиномиального автоморфизма:
его нули являются конкретными исключительными слоями, требующими
проверки.

\subsection{Строгая обструкция одного естественного переноса для $H(3,7,8)$}
Тождество
\[
 C_2(C_3)C_8-C_7^2+2=C_2
\]
даёт модификацию
\[
 \beta(x,y,z)=(x,(x^2-2)z-y^2+2,y),
\]
которая переводит $(C_3,C_7,C_8)$ в $(C_3,C_2,C_7)$.
Последняя тройка допускает относительную координату
$H_0(x,Y,Z)=xY^2-2x-Z$, поскольку
$C_3(C_2^2-2)-C_7=t$. Однако
\[
 H_0(x,2-y^2,y)=x(2-y^2)^2-2x-y
\]
имеет степень 4 по $y$ при $x^2=2$.

Более сильно, замена $H_0$ на другой многочлен $H$, удовлетворяющий
$H(C_3,C_2,C_7)=t$, не может сделать центр линейным на обоих
исключительных слоях. Действительно, пусть $r^2=2$ и предположим
$H(r,2-y^2,y)=\alpha y+\beta$.
В $k[t]/(t^3-3t-r)$ имеем
\[
 C_7=t-r(t^2-2),\qquad 2-C_7^2=C_2.
\]
Следовательно, $t=\alpha(t-rt^2+2r)+\beta$.
Базис $1,t,t^2$ даёт одновременно $\alpha=1$ и $\alpha r=0$,
что невозможно. Это доказывает невозможность применения именно
линейного критерия центра к этой модификации и этой параметризации.
Это \textbf{не} доказательство невыпрямляемости $H(3,7,8)$.

\subsection{Нижняя граница для относительной координаты $H(3,7,8)$}
Все многочлены общей степени не выше 4, удовлетворяющие
$H(C_3,C_7,C_8)=t$, образуют двумерное аффинное семейство
$H=H_0+\lambda F_0+\mu F_1$, где
\begin{align*}
H_0={}&x^3z+x^3-x^2yz-xy^2-2xz-x+y^3+2yz-2y,\\
F_0={}&x^3y-x^2z-3xy+y^2-z^2+2z,\\
F_1={}&x^4-x^2z^2+x^2z-3x^2+xy^3-xyz-2xy\\
 &+y^2z-y^2-z^3+2z^2-z+2.
\end{align*}
Это полное решение конечной линейной системы по всем мономам общей
степени не выше 4. Для любого $\lambda,\mu$ справедливо
\[
 H_y(0,0,1)=H_z(0,0,1)=0,\qquad H_x(0,0,1)=-3.
\]
Поэтому $(x,H)$ нельзя дополнить до полиномиального автоморфизма,
сохраняющего $x$. Следовательно, относительная координата параметра,
если существует, должна иметь общую степень как минимум 5.
Это утверждение касается только автоморфизмов, сохраняющих $x=C_3$;
оно не исключает произвольного выпрямления с параметрической
компонентой степени 4.

\subsection{Ограниченный квадратичный поиск}
Для координат $(x,y,z)=(C_3,C_8,C_7)$ положим
$k=x^2-1$, $\Delta=k^2-x^2$ и $L=z+A(x)y$.
Тогда $L=ut+vs$, где $u=k-Ax$, $v=-x+Ak$.
Рассматривался анзац
\[
 H=P(x)L^2+Q(x)L+R(x)+M(x)y,
 \qquad H(C_3,C_8,C_7)=t,
\]
с $M=mx^e$. Сравнение коэффициентов в базисе $1,t,s$ даёт
необходимое тождество
\[
 v+P\bigl(u^3-3uv^2-xv^3\bigr)+mx^e\Delta=0.
\]
Его коэффициенты проверялись методом базисов Грёбнера. Для
$A=cx^j$, $j=0,1,2,3$, $1\leq e\leq16$, при неотрицательной
ожидаемой степени $P\leq6$ все исследованные идеалы равны $(1)$.
Отдельно проверены нечётные $A$ степеней не выше 1, 3, 5 и нечётные
$P$ степеней не выше 1, 3 (а также случаи $(1,5),(3,5)$):
во всех этих системах также получен базис $[1]$.
Степень $e$ во втором переборе равна $3\deg A+\deg P+3$.
Это точный отрицательный результат для выписанных конечных систем.
Он не исключает произвольных $A,P,M$, иных степеней, иных относительных
координат или конструкций с сохранением другой первой координаты.

\subsection{Полином параметра для $H(4,5,7)$}
Точная линейная алгебра по мономам показала, что минимальная общая
степень полинома $F$ с $F(C_4,C_5,C_7)=t$ равна 3, причём в степени
не выше 3 решение единственно:
\[
 F(x,y,z)=(x^2-2)z-y^3+3y.
\]
Этот конкретный $F$ не является координатой: его градиент обращается
в нуль при $x^2=2$, $y^2=1$, $z=0$. Следовательно, любое выпрямление
данной параметризации должно использовать полином параметра общей
степени как минимум 4 (или после аффинной перепараметризации ---
соответствующий аффинный образ полинома степени как минимум 4).
Это не препятствие существованию такого полинома большей степени.

\section{Что остаётся открытым и как устроен следующий поиск}

\paragraph{Доказанные границы.}
Выпрямление всех гармонических вложений не получено. В частности,
$H(3,7,8)$, $H(3,7,11)$ и $H(4,5,7)$ в этом досье не имеют
ни удостоверенного выпрямителя, ни доказательства невыпрямляемости.
Невозможность одного переноса, наличие критических точек у одного
полинома параметра и неудача ограниченной системы Грёбнера не заменяют
ни одного из этих двух результатов. Для утверждений о неприводимости
случая нужно исключить все допустимые координаты, а не один анзац.

\paragraph{Конкретная последовательность дальнейших задач.}
\begin{enumerate}
\item Для заданного $H(a,b,c)$ сначала применять точные чебышёвские
инволюции и сдвиги индексов; сохранять сам автоморфизм и обратный,
а не только конечную тройку индексов.
\item При кубической первой координате использовать матрицу остатков
в базисе $1,t,t^2$ и доказанный фильтр допустимых исключительных слоёв.
Для переноса с сохранённой второй координатой $C_7$ возможный
множитель имеет корни только в $x=0$. Естественный переход через
$x^2-2$ поэтому отбрасывается до большого символьного раскрытия.
\item Искать другой полином параметра в классе $R+I(\gamma)$.
Для кандидата проверять геометрические нули соответствующего
якобиева LND и строить слайс точным обращением по модулю полного
исключительного множителя. Отсутствие критических точек у кандидата
само по себе не удостоверяет координатность.
\item Для $H(4,5,7)$ начинать поиск координаты параметра со степени 4:
единственный полином степеней до 3 уже исключён. Для каждого найденного
полинома нужны компаньоны и полиномиальная обратная замена.
\item Для связанного кольца $B=\ker D_\gamma$ изучать существование
LND, перемещающего $t$. Дихотомия $\mathrm{ML}(B)\in\{k,k[t]\}$
уточняет цель, но не предоставляет алгоритма перечисления всех LND.
\end{enumerate}

\paragraph{Какие следствия уже можно использовать.}
Для удостоверенных моделей известны: два образующих идеала кривой
в координатах выпрямителя; свободное действие $\mathbb G_a$, одна
из орбит которого есть эта кривая; полиномиальные ядра связанных
пятимерных LND; явная тривиализация указанных конкретных действий;
минимальность степени числового поля определения, равной двум,
для нетривиальных вещественных базовых узлов. Ни одно из этих
утверждений не требует мотивных или гомологических методов.

\section{Материалы и повторение вычислений}

Архив содержит это досье с исходниками LaTeX, новые вычислительные
скрипты, JSON-сертификаты, журналы и манифест исходной версии репозитория.
Внешний репозиторий указан как фиксированная зависимость; его
полную копию архив не дублирует. Файл \texttt{bootstrap\_source.py}
получает только нужные файлы указанного коммита и проверяет Git-хеш
каждого файла. Это нужно лишь проверкам, зависящим от базового выпрямителя.
Самостоятельные сертификаты минимальных степеней слайсов используют
стандартную библиотеку Python и не требуют загрузки репозитория.

\begin{verbatim}
python -m pip install -r requirements.txt
python bootstrap_source.py
python research/repo_audit/independent_base_lnd.py
python research/harmonic_literature/periodic_shears.py
python research/lnd_certificate/verify_lnd.py
python research/lnd_certificate/verify_minimal_slice.py
python research/repo_audit/review_minimal_slice.py
python research/harmonic_compute/verify_new_results.py
python research/harmonic_compute/reduction_scan.py
\end{verbatim}
Подробный перечень команд дан в \texttt{README.txt}. Результаты
исходного запуска находятся рядом со скриптами. Новые проверки
используют явные исключения или явные проверки результата, где
это предусмотрено кодом; запускать команды следует обычным Python,
без отключения стандартных проверок флагом \texttt{-O}.

\paragraph{Основные первоисточники и зачем их читать.}
\begin{enumerate}
\item G.~Freudenburg, J.~Freudenburg,
\href{https://arxiv.org/abs/0902.3440}{Curves Defined by Chebyshev Polynomials}
(2009), предложения 3.1, 4.1, 4.2. Критерий комплексного вложения,
узлы плоской проекции и алгебраическая редукция индексов.
\item P.-V.~Koseleff, D.~Pecker,
\href{https://arxiv.org/abs/0909.3281}{Chebyshev diagrams for two-bridge knots};
\href{https://arxiv.org/abs/1203.4376}{Harmonic knots}.
Топология гармонических узлов, в частности классификация при $a\leq4$;
этот результат следует отделять от комплексной полиномиальной эквивалентности.
\item A.~van den Essen, P.~van Rossum,
\href{https://doi.org/10.1090/S0002-9939-01-06212-8}{Triangular derivations related to problems on affine $n$-space}
(2002), Theorem 4.1, Conjectures 5.1--5.3. Точное место гипотезы E
и связь между выпрямлением, ядрами LND и конкретным кандидатом
к проблеме сокращения.
\item D.~Lewis,
\href{https://arxiv.org/abs/2203.08178}{A unified approach to embeddings of a line in 3-space}
(2022), Lemma 1, Theorem 2, Corollary 4. Подъём параметра до
координаты и устранение полюсов рациональных преобразований.
\item K.~Masuda,
\href{https://doi.org/10.1016/j.jalgebra.2008.12.012}{Homogeneous locally nilpotent derivations having slices and embeddings of affine spaces}
(2009), и
\href{https://sci-tech.ksc.kwansei.ac.jp/~kayo/cmn/papers/eembed.pdf}{Equivariant embeddings and semi-invariant locally nilpotent derivations}.
Связь с градуировками и симметриями; дополнительные условия на веса
нужны для обратных импликаций.
\end{enumerate}
Нумерация и формулировки сверялись с текстами источников, а не только
с поисковыми аннотациями. Даты, автоматически вставленные в HTML
старых arXiv-статей, не использовались как даты публикации.

\end{document}
